% ============================================================
\section{Randoop}
\label{sec:randoop}
% ============================================================

Randoop genera sequenze casuali di chiamate ai metodi pubblici della classe indicata e le esegue,
classificando ogni sequenza come test di regressione (nessuna eccezione, il comportamento
osservato viene cristallizzato in un'asserzione) o come test error-revealing (un'eccezione non
attesa, sintomo di una violazione di contratto Java).

\subsection{BrokerImpl}

Il jar di Randoop è stato scaricato da GitHub (versione 4.3.3).

È stato compilato \texttt{openjpa-kernel} per avere i riferimenti su cui lavora Randoop:

\begin{lstlisting}
mvn -pl openjpa-kernel -am compile
\end{lstlisting}

È stato estratto il classpath delle dipendenze del modulo, necessario per eseguire Randoop
insieme al codice appena compilato:

\begin{lstlisting}
mvn -pl openjpa-kernel dependency:build-classpath \
  -Dmdep.outputFile=kernel-cp.txt
\end{lstlisting}

Il classpath da passare a Randoop è la concatenazione di \texttt{openjpa-kernel/target/classes}
con il contenuto di \texttt{kernel-cp.txt}.

È stato eseguito un primo tentativo esplorativo, con un limite di tempo di 30 secondi e nessun
limite sul numero di test, per capire quanti test genera Randoop senza vincoli sull'output: il
risultato è stato quasi 1800 sequenze. Randoop nasce per essere uno strumento automatico, quindi
di per sé non richiede una revisione manuale dell'output generato; in questo contesto ipotizzo però
ragionevole voler osservare concretamente i risultati
prodotti, cosa che con quasi 1800 sequenze diventa impraticabile. Per questo la generazione
definitiva usa \texttt{--output-limit=100}.

È stata lanciata la generazione definitiva su \texttt{BrokerImpl}:

\begin{lstlisting}
java -Xmx3000m -cp tools/randoop-all-4.3.3.jar:openjpa-kernel/target/classes:<contenuto-di-kernel-cp.txt> \
  randoop.main.Main gentests \
  --testclass=org.apache.openjpa.kernel.BrokerImpl \
  --output-limit=100 \
  --junit-output-dir=report/data/randoop-brokerimpl
\end{lstlisting}

L'esecuzione produce 101 sequenze in totale (162 tentativi):
12 classificate error-revealing, scritte in \texttt{ErrorTest0.java}; delle restanti, 88
candidate come regressione, 61 diventano il file finale \texttt{RegressionTest0.java} dopo la
selezione interna dello strumento. L'output completo stampato da Randoop durante questa
esecuzione è salvato in \texttt{report/data/randoop-brokerimpl/randoop-log.txt}.

I due file generati si trovano fuori dal ciclo maven, quindi Maven/Surefire non
li vede. Sono stati compilati ed eseguiti con gli
strumenti del JDK direttamente:

\begin{lstlisting}
javac -cp openjpa-kernel/target/classes:<contenuto-di-kernel-cp.txt>:junit-4.13.2.jar \
  -d <dir-output> RegressionTest0.java ErrorTest0.java

java -cp <dir-output>:openjpa-kernel/target/classes:<contenuto-di-kernel-cp.txt>:junit-4.13.2.jar:hamcrest-core.jar \
  org.junit.runner.JUnitCore RegressionTest0

java -cp <dir-output>:openjpa-kernel/target/classes:<contenuto-di-kernel-cp.txt>:junit-4.13.2.jar:hamcrest-core.jar \
  org.junit.runner.JUnitCore ErrorTest0
\end{lstlisting}

I 61 test di regressione passano tutti. I 12 error-revealing falliscono tutti con
\texttt{NullPointerException}, con la stessa causa in ognuno: \texttt{new BrokerImpl()} lascia i
campi \texttt{\_conf}, \texttt{\_transEventManager}, \texttt{\_lifeEventManager},
\texttt{\_cache} e \texttt{\_log} a \texttt{null}, impostati solo da \texttt{initialize(...)}.

\subsection{LoginDialog}

